\section{Discussion}\label{sec:discussion}

\subsection{Mechanisms}

The evidence describes a minimum wage that moved workers rather than wages. The reallocation is my most robust result: formality falls and self-employment rises, the shifts scale with the regional bite of the floor, and they are confirmed by a continuous-exposure design under randomization inference. The natural mechanism is the classic one, a binding floor in the covered sector inducing reallocation to the uncovered sector, the displacement channel that \citet{jaramillo_lopez_2006} documents for Peru and \citet{ham_2018_honduras} for Honduras. Its incidence is exactly where that mechanism predicts: among the young, among women, and in high-bite departments, the workers and places where the new floor cut deepest into the wage distribution. The panel transitions of Section~\ref{sec:transitions} identify the flow behind the stock change and refine the mechanism in one important way: retention of incumbent formal workers did not fall, while entry into formal jobs from informality dropped by nearly half, so the reallocation operated through the creation of formal matches rather than through their destruction. The floor did not push incumbents out; it closed the door on new formal matches for low-educated workers. This is the same hiring-margin adjustment that \citet{dustmann_2022} document for Germany, and it explains why the burden falls on marginal entrants, secondary earners, the young, and women, rather than on protected incumbents. The accompanying rise in hours is consistent with movement toward self-employment, where hours are less constrained. One nuance qualifies the destination of the displaced margin: in the regional design the unconditional self-employment rate does not rise with exposure (Appendix Table~\ref{tab:uncond}), so across departments the robust absolute margin is the destruction of formal employment, and the rise in the self-employment share of employment partly reflects the shrinking employment base; the absolute shift into own-account work appears in the education contrast, where the displaced workers and their closest substitutes are compared directly.

Why did no aggregate pay response accompany the reallocation? Not because the floor is a dead letter: the first-stage evidence shows that within the formal sector the spike at the old minimum migrated almost completely to the new one, so covered employers complied. The explanation is the arithmetic of coverage in a two-sector labor market \citep{ulyssea_2018}. Only 13 percent of employed low-skilled workers are formal, so even full compliance inside the formal segment moves the group's average wage very little; meanwhile pay in the informal segment drifted \emph{further} below the floor, and the two responses cancel in the aggregate. Inflation compounds this: the 2022 increase of ten percent in nominal terms met consumer price inflation of roughly eight percent, so the real value of the floor barely rose and then eroded. A floor that is honored only where labor costs are enforceable cannot lift the group's wage distribution, but it can still shape which employment arrangements firms and workers choose, and on my estimates it did. This is the same reallocation logic that \citet{dustmann_2022} document across establishments in Germany, transposed to the margin that matters in a developing economy: the boundary between formal and informal work \citep{jales_2018}.

Two caveats discipline this reading. The reallocation, while robust to specification and to design-based inference, is sensitive to modest violations of parallel trends in the education contrast \citep{rambachan_roth_2023}, and it does not reappear around the 2025 reform, where the same design and the same magnitudes yield nulls on every margin and a small opposite-signed self-employment estimate. I interpret the contrast between the two episodes as state dependence, and Section~\ref{sec:reconciliation} probes the two candidate state variables directly: the evidence favors the position of the floor in the formal pay distribution, the 2022 increase swept a band holding a quarter of low-educated formal pay while the 2025 increase swept a thinner one, and it rejects the labor-market-churn version, since formal match creation was in fact higher on the eve of 2025. The entry-margin evidence of Section~\ref{sec:transitions} identifies the margin through which any such state dependence must operate: the creation of new formal matches, whose profitability the floor governs. I cannot, however, fully exclude the alternative reading, that residual recovery dynamics correlated with education contribute to the 2022 compositional estimates; the dose-response and continuous-exposure evidence is what tilts me toward a causal interpretation, since differential recovery by education has no reason to track the pre-reform regional Kaitz index.

\subsection{Reconciliation with the previous literature}\label{sec:reconciliation}

My results sit comfortably with some of the prior evidence and in tension with the rest, and the pattern is itself informative. They echo the older Peruvian work of \citet{jaramillo_lopez_2006}, who found the 2003 reform largely innocuous with displacement toward informality, and of \citet{cespedes_2005}, who found effects concentrated among the young. They are consistent with the broad international finding that minimum wage increases in this range produce small employment effects \citep{card_krueger_1995, cengiz_2019}.

The tension is with the most recent Peruvian study. \citet{jimenez_2023} finds that the 2016 reform raised formal wages by four to five percent and increased formality, a lighthouse pattern that I do not reproduce for 2022. Three differences plausibly account for the divergence. The 2016 increase occurred in a low-inflation period, so its real bite was larger and more durable than that of the inflation-eroded 2022 increase. The identification differs: \citet{jimenez_2023} uses a department-level dose defined on the formal wage distribution, while I contrast low- and high-education workers at the individual level, a design that places more weight on the large informal segment of the low-skilled. And the macroeconomic setting differs, stable growth in 2016 against a post-pandemic recovery in 2022. One observation helps to sort out these explanations: 2016, the lighthouse episode, was itself a calm year, and so was 2025, in which I find no effects at all, so macroeconomic state alone cannot be the whole story. What separates 2016 from 2025 is the \emph{level} of the floor when the increase arrived: in 2016 the minimum stood well below the median wage and a real increase could plausibly propagate as a norm, whereas by 2025 the floor already exceeded the national median, so a further increase started from a level that practice had already ceased to honor for the marginal worker. On this reading the three episodes are ordered by two state variables, namely how high the floor already sits in the wage distribution and how much re-sorting the labor market is doing: a moderate floor rising in a stable market can lift norms (2016); a high floor rising amid post-pandemic re-sorting moves the composition margin (2022); and a high floor rising in a settled market moves nothing (2025).

Because a story with two free state variables can rationalize almost anything, I probe both directly and let the data discard one. The floor-position variable behaves as the story requires: on the eve of the 2022 reform, 24 percent of low-educated formal wage earners were paid inside the band the incoming floor was about to sweep and 16 percent bunched at the old floor, whereas on the eve of the 2025 reform those shares had fallen to 17 and 14 percent. The re-sorting variable does not: the new-hire share among low-educated formal workers was \emph{higher} in 2024 (41 percent) than before the 2022 reform (32 percent), so the 2025 null cannot be attributed to a scarcity of new formal matches. A department-level dose test points the same way: interacting the education contrast with pre-reform formal match creation yields no gradient in either episode, with or without the Kaitz control.\footnote{The state-variable moments and the dose estimates are included in the replication package (\texttt{state\_premise.csv}, \texttt{state\_dose.csv}).} I therefore retire the churn version of the argument and rest the interpretation on the floor's position: the entry margin closes when the incoming floor sweeps a dense band of formal pay, as in 2022, and has little left to close when the floor has already moved past the range that formal pay practice honors, as in 2025. The swept-band contrast is real but modest, so I present this as the reading most consistent with the evidence rather than as established.

\subsection{External validity and limitations}

Five limitations bound the interpretation. First, identification rests on the assumption that, absent the reform, low- and high-education workers would have followed parallel paths. I test this where possible and bound the consequences of its failure; the formality pre-trends are flat, but the wage pre-trend test rejects on account of a single early-2021 quarter, and the employment series carries a clear recovery pre-trend, which is why I state the wage conclusion as ``no detectable gain'' rather than a precise null and treat the employment-level result as inconclusive. Second, the design identifies effects on low-skilled workers relative to high-skilled workers. The comparison group is not entirely unexposed (40 percent of high-skilled wage earners sat below the new floor pre-reform), so the estimates are relative objects that, to first order, understate the effect on the fully exposed; the continuous-exposure design partly addresses this by dispensing with the comparison group. Third, the main design uses repeated cross-sections; the ENAHO Panel analysis of Section~\ref{sec:transitions} addresses the flow-versus-stock distinction directly, but it is annual rather than quarterly, its pre-reform benchmark includes transitions out of the pandemic year 2020, and its formal-base samples are modest, so I treat it as corroborating the mechanism rather than as delivering precise magnitudes. Fourth, informality in 2024 is measured with a reconstructed indicator rather than the official one; within the 2025 validation window the indicator is reconstructed consistently in every quarter, but the 2024 splice adds measurement error to the tail of the 2022 event study. Finally, the window spans at most ten clean post-reform quarters, so I speak to short- and medium-run effects and not to the long-run adjustment of capital or firm entry and exit.

\subsection{Policy implications}

The results carry a sober message for minimum wage policy in a high-informality economy. A floor that goes unenforced for most of its intended beneficiaries does little to raise their pay, and where it does bind, in the formal sector, it may nudge some workers toward less protected arrangements. Raising the minimum wage is therefore an imperfect instrument for lifting the earnings of the low-skilled in Peru, a conclusion that echoes \citet{jaramillo_lopez_2006}. This is not to say the minimum wage is harmless or useless; my estimates are relative and my window is short. But it does imply that formalization and enforcement are preconditions for the floor to reach the workers it is designed to protect. Absent them, the margin of adjustment is the composition of employment rather than the level of pay.
